\chapter{Le pitch, et ce qu'on fait ce soir}
\label{ch:pitch}

\section{Les quinze minutes}

Trois voix, une par module : le jury entend chaque responsable parler de ce qu'il a construit.

\begin{center}\small
\begin{tabularx}{\textwidth}{@{}l L{2.6cm} Y L{2.3cm}@{}}
\toprule
\textbf{Temps} & \textbf{Partie} & \textbf{Contenu} & \textbf{Qui} \\
\midrule
0:00 -- 1:30 & Le problème & Un client écrit à sa banque ; un Français sur quatre est client de
  La Banque Postale ; l'expérience client est sa priorité n\textsuperscript{o}\,1 ; la règle :
  open source, local, frugal. & Yahya \\
1:30 -- 4:00 & La démarche & Les quatre temps ; le constat sur Banking77 ; la cascade (une
  figure) ; le gros modèle à l'instant précis. & Yahya \\
4:00 -- 10:00 & La démonstration & Le script du chapitre~\ref{ch:demo}. & Yahya pilote,
  Sohayb commente \\
10:00 -- 12:30 & Les résultats & Compréhension sur les trois jeux de test ; scénarios ; tableau
  A/B/C ; coût de la préparation hors ligne. & Younes \\
12:30 -- 14:00 & L'explicabilité & Raisons et produits écartés ; garde-fou ; règlement européen
  sur l'IA ; mention légale. & Sohayb \\
14:00 -- 15:00 & Limites et suite & Ce qui est simulé ; ce qu'il faudrait pour 17\,000 points de
  contact. & Younes \\
\bottomrule
\end{tabularx}
\end{center}

\section{Trois messages}

\begin{enumerate}
  \item \og Notre assistant comprend souvent avec de simples mots-clés, sinon avec 118 millions de
        paramètres ; il décide avec des règles que le client peut lire ; il n'appelle un LLM
        qu'une fois, à la fin, pour mettre en mots ce qui est déjà calculé. \fg
  \item \og Nous l'avons mesuré contre une version où un LLM de 14 milliards de paramètres fait
        tout. \fg{} (suivi des chiffres)
  \item \og Chaque conseil dit pourquoi --- et pourquoi pas. \fg
\end{enumerate}

\section{Les questions probables du jury}

\begin{center}\small
\begin{longtable}{@{}L{4.6cm} L{9.8cm}@{}}
\toprule
\textbf{Question} & \textbf{Réponse préparée} \\
\midrule
\endhead
Pourquoi ne pas tout confier à un LLM ? & Parce qu'on l'a mesuré : la configuration A fait tout
  avec un LLM de 14 milliards de paramètres ; voici sa qualité, son énergie et sa latence face à
  la cascade. \\
\addlinespace
Comment mesurez-vous l'énergie ? & CodeCarbon en mode hors ligne, par lots de 200 requêtes,
  puissance de repos soustraite, trois répétitions. Une requête seule est trop courte pour être
  mesurée honnêtement. \\
\addlinespace
Et si le classifieur se trompe ? & En dessous du seuil $\tau$, il ne devine pas : il pose une
  question fermée, puis escalade. $\tau$ a été choisi par balayage pour 97\,\% de précision sur
  les messages acceptés. Voici la courbe. \\
\addlinespace
Le LLM peut-il inventer un taux ? & Il ne reçoit que des faits calculés, et chaque nombre de sa
  réponse est contrôlé contre ces faits ; en cas d'écart, la réponse gabarit s'affiche. Voici le
  taux d'échec du garde-fou. \\
\addlinespace
Pourquoi Ministral ? & Parce qu'en descendant depuis le 14B, c'est le plus petit qui tient notre
  seuil de qualité en français --- voici la courbe --- et qu'il est européen. Si un plus petit
  avait suffi, le protocole l'aurait retenu. \\
\addlinespace
Des mots-clés, en 2026 ? & Oui : ils traitent une part mesurée des demandes pour un coût quasi nul,
  et ils disent quand ils ne sont pas sûrs. Le reste passe au classifieur. Moins de LLM, plus
  d'algorithmes. \\
\addlinespace
Bank Marketing, ce sont des données portugaises de 2008 ? & Oui. Elles ancrent des profils
  réalistes et donnent un signal d'appétence faible, présenté comme tel. L'adaptateur permet de
  brancher une vraie base sans toucher au reste. \\
\addlinespace
Et la conformité ? & Données simulées et traitement local. Le règlement européen classe à haut
  risque la notation de solvabilité : nous n'en faisons pas, un conseiller décide. Transparence :
  l'utilisateur sait qu'il parle à une IA. Mention légale du crédit. \\
\addlinespace
Comment passer à 17\,000 bureaux ? & Le coût par requête est mesuré, donc extrapolable. Les étages
  1 à 4 tournent sur CPU ; seul le petit LLM a besoin d'un GPU, et un serveur partagé suffit. Les
  pistes : voix locale et réponse dans la langue du client, dans la ligne de l'outil multilingue
  annoncé pour novembre. \\
\addlinespace
Pourquoi pas d'affinage ? & Deux jours, des GPU partagés, et une tête de classification entraînée
  en secondes qui suffit. C'est aussi la conclusion d'un système en production que nous avons
  construit : la recherche bat l'affinage quand le modèle de base évolue vite. \\
\addlinespace
Que se passe-t-il hors périmètre ? & L'assistant le dit, et propose un conseiller. Il ne fait pas
  semblant. \\
\bottomrule
\end{longtable}
\end{center}

\section{Ce qu'on fait ce soir}

\begin{enumerate}
  \item Lire la page de synthèse et le chapitre~\ref{ch:organisation} (vingt minutes) ; valider ou
        amender la répartition, puis ne plus y revenir.
  \item Écrire \f{src/contrats.py} ensemble (trente minutes).
  \item Rendre le dépôt GitHub privé : il est public, et il va contenir des éléments tirés de
        supports marqués \og Interne \fg.
  \item Déplacer les données téléchargées dans \f{data/raw/} et l'ajouter au \f{.gitignore}.
  \item \textbf{Younes} : écrire à Denis Barthou pour DeepEirb et le modèle vLLM ; lancer la
        traduction, ou décider du repli.
  \item \textbf{Yahya} : squelette de \f{app.py} et de \f{reply()} avec bouchons ; parcours carte.
  \item \textbf{Sohayb} : les 8 clients et le catalogue v0 ; squelette du rejeu des scénarios.
  \item En filaire, télécharger \f{ministral-3:3b}, \f{qwen3:1.7b}, \f{gemma3:1b} et
        \f{multilingual-e5-small} sur un poste de la salle.
  \item Vérifier J0 avant de partir : \f{streamlit run app.py} répond.
\end{enumerate}
